\documentclass[aspectratio=169]{beamer}
\usetheme{metropolis}
\usepackage{nexus_theme}

\title{Research Methodologies}
\subtitle{Module 0.2: Quantitative, Qualitative, and Mixed Methods}
\author{Nexus Scholar Suite --- Modern Research Methodologies}
\date{}

\begin{document}

% --- TITLE SLIDE ---
\begin{frame}[plain]
  \titlepage
\end{frame}

% --- LEARNING OBJECTIVES ---
\begin{frame}{What You'll Learn}
  By the end of this lesson, you will be able to:
  \begin{enumerate}
    \item \textbf{Define} quantitative, qualitative, and mixed-methods paradigms.
    \item \textbf{Identify} appropriate methodologies for different questions.
    \item \textbf{Evaluate} the strengths and limitations of each approach.
    \item \textbf{Explain} validity and reliability in research design.
  \end{enumerate}
\end{frame}

% --- THE THREE PARADIGMS ---
\begin{frame}{The Three Paradigms}
  \begin{columns}[T]
    \begin{column}{0.33\textwidth}
      \begin{block}{Quantitative}
        \begin{itemize}
          \item Numbers \& Stats
          \item Objective
          \item Deductive (Test theory)
          \item e.g., Surveys, RCTs
        \end{itemize}
      \end{block}
    \end{column}
    \begin{column}{0.33\textwidth}
      \begin{exampleblock}{Qualitative}
        \begin{itemize}
          \item Words \& Meaning
          \item Subjective
          \item Inductive (Build theory)
          \item e.g., Interviews
        \end{itemize}
      \end{exampleblock}
    \end{column}
    \begin{column}{0.33\textwidth}
      \begin{alertblock}{Mixed Methods}
        \begin{itemize}
          \item Both combined
          \item Triangulation
          \item Complete ecosystem view
        \end{itemize}
      \end{alertblock}
    \end{column}
  \end{columns}
\end{frame}

% --- THE FOREST ANALOGY ---
\begin{frame}{The Forest Analogy}
  \begin{center}
  \resizebox{0.92\linewidth}{!}{
  \begin{tikzpicture}[node distance=1.2cm and 1.8cm]
    \node[card, fill=NexusBlue!10, draw=NexusBlue, text width=3.5cm] (quant) {
      \textbf{Satellite View}\\
      (Quantitative)\\
      \textit{Total biomass, broad trends, high scale, low context.}};
      
    \node[successcard, right=of quant, text width=3.5cm] (qual) {
      \textbf{Ground View}\\
      (Qualitative)\\
      \textit{Studying one tree's root system. Deep context, low scale.}};
  \end{tikzpicture}
  }
  \end{center}
\end{frame}

% --- MATCHING METHOD TO QUESTION ---
\begin{frame}{Matching Methodology to the Question}
  \begin{table}[]
    \begin{tabular}{|l|l|}
    \hline
    \textbf{Research Question Type} & \textbf{Paradigm} \\ \hline
    Does X cause Y? & \textbf{Quantitative} (Experimental) \\ \hline
    How many people do Z? & \textbf{Quantitative} (Descriptive) \\ \hline
    How do people experience X? & \textbf{Qualitative} (Phenomenology) \\ \hline
    Why does culture Y behave this way? & \textbf{Qualitative} (Ethnography) \\ \hline
    Does X work, and why do people like it? & \textbf{Mixed Methods} \\ \hline
    \end{tabular}
  \end{table}
  \vspace{0.3cm}
  \textit{Rule: Never pick your tool before you define your question.}
\end{frame}

% --- VALIDITY AND RELIABILITY ---
\begin{frame}{Validity vs. Reliability}
  \begin{block}{Two distinct requirements for good science}
    \begin{itemize}
      \item \textbf{Reliability (Consistency):} If someone else repeats the exact same study, will they get the exact same result?
      \item \textbf{Validity (Accuracy):} Are you actually measuring what you claim to be measuring?
    \end{itemize}
  \end{block}
  \vspace{0.3cm}
  \textbf{Analogy:} A broken scale that always reads 100 lbs is perfectly \textit{reliable}, but completely \textit{invalid}.
\end{frame}

% --- DATA SATURATION ---
\begin{frame}{Sample Sizes: Power vs. Saturation}
  \begin{columns}[T]
    \begin{column}{0.48\textwidth}
      \begin{block}{Quantitative: Statistical Power}
        Driven by math. You need enough $N$ to prove that your result isn't just random chance (p < 0.05). Often requires hundreds or thousands of subjects.
      \end{block}
    \end{column}
    \begin{column}{0.48\textwidth}
      \begin{exampleblock}{Qualitative: Data Saturation}
        Driven by information. You stop interviewing people when the newest interview yields no new insights or themes. Often reached at $N=12$ to $N=20$.
      \end{exampleblock}
    \end{column}
  \end{columns}
\end{frame}

% --- KEY TAKEAWAY ---
\begin{frame}{Key Takeaway}
  \begin{block}{What We Accomplished}
    \begin{enumerate}
      \item \textbf{Paradigms}: Quant (Numbers) vs Qual (Words).
      \item \textbf{Alignment}: The research question dictates the methodology.
      \item \textbf{Metrics}: Reliability (Consistency) vs Validity (Accuracy).
    \end{enumerate}
  \end{block}
  \vspace{0.3cm}
  \textbf{Next:} Module 0.3 --- The Publishing Ecosystem
\end{frame}

% --- EXERCISE PREVIEW ---
\begin{frame}{Your Turn: Exercise}
  \begin{exampleblock}{Exercise: Methodological Matchmaker}
    Review three flawed research proposals, identify the methodological mismatch, and redesign the approach to correctly answer the core research question.
  \end{exampleblock}
\end{frame}

% --- STANDOUT CLOSE ---
\begin{frame}[standout]
  A perfect statistical analysis of the wrong question\\
  is still the wrong answer.
\end{frame}

\end{document}
